\section{刑罚的裁量}
\begin{enumerate}
  \item \zx{量刑情节}。从重、从轻、减轻、免除处罚。
        \begin{enumerate}
          \item 从重、从轻区别。同一个刑格 3 - 7 年。
                \begin{enumerate}
                  \item 加重/法定刑升格 \yjt \dy 7 年。上一个台阶。
                  \item 从重 \yjt 7 年。
                  \item 从轻 \yjt 3 年。\xjt \zx{未遂、立功、立功}。
                  \item 减轻 \yjt \xyh 3 年。下一个台阶。
                \end{enumerate}
        \end{enumerate}
  \item \zx{累犯}。前罪满 18 不要忘。
        \begin{enumerate}
          \item 一般累犯。
                \begin{enumerate}
                  \item 前后都是故意犯罪 + 有期。
                  \item 后罪发生在前罪刑罚执行完毕（只看主刑）或赦免 5 年内。
                  \item 前罪 + 成功的缓刑/假释 \yjt 不构成/有可能构成。
                \end{enumerate}
          \item \zd{特殊累犯。前后都是国恐黑 3 种}。
          \item 应当从重处罚。不得适用缓刑、假释。
        \end{enumerate}
  \item 自首。
        \begin{enumerate}
          \item 一般自首。\bs{一审判决前 + 自动投案 + 如实供述}。如实供述不要忘！！
                \begin{enumerate}
                  \item 犯罪后，被抓前。最迟一审判决前。
                  \item 自动投案。\zd{xx 行为进行投案，属于自动投案}。\al{省略 1W 字，……，逃离现场 \cuo}。
                        \begin{enumerate}
                          \item 自动 = 不是被抓。\zd{尚未受到讯问、未被采取强制措施时}。
                          \item 投案 = 主动置于司法机关控制之下。
                        \end{enumerate}
                  \item 如实供述。\al{甲买水果与人争执，水果刀刺中乙后背。甲坐在原地等待警察到来。甲的行为属于自首 \cuo{甲仅仅只是自动投案而已}}。
                        \begin{enumerate}
                          \item \zd{如实。一五一十}、共犯供全部。反面 \yjt 否认自己的行为。
                          \item 供述。主要犯罪事实。只关注罪，不关注证。\al{拒绝说明[\xhx{凶器藏匿地点}、真实身份] \dui}。
                          \item vs 辩解。\al{主张自己是正当防卫 \dui}。
                        \end{enumerate}
                  \item \zd{各种奇葩情况}。
                        \begin{enumerate}
                          \item \al{抓后逃 \yjt 不构成}。
                          \item \al{主动投案 + 如实供述 \yjt 逃 \yjt 再自首 \yjt 构成}。
                          \item \al{A、B、C 只供了 A、B \yjt A、B 构成自首，C 不构成}。
                        \end{enumerate}
                \end{enumerate}
          \item \zx{特别自首/准自首/余罪自首}。\zd{未被掌握}、证据不足、能成立。\zd{主动交代了另一起犯罪事实，构成准自首}。\al{有证据，补充完善 \cuo}。
          \item 单位自首。
                \begin{enumerate}
                  \item 主管自首，视为单位自首。
                  \item 直接责任人自首，只能算个人自首。
                  \item 单位自首，主管、责任人被动归案 + 坦白，可以算自首。大树底下好乘凉。
                \end{enumerate}
          \item \zx{vs 坦白}。主动交代的罪行\zd{已为司法机关掌握}。\zd{如实供述 xx 事实，属于坦白。否认 xx，对此不能认定为坦白}。
        \end{enumerate}
  \item \bs{立功}。\zx{揭发他人、协助抓捕}。亲友实施不构成。
        \begin{enumerate}
          \item 揭发他人犯罪。表现为\zd{提供合法得来的犯罪线索、新信息}。
                \begin{enumerate}
                  \item 合法（扩大解释）。反面排除。\zd{贿买、暴力、胁迫}等非法手段。
                        \begin{enumerate}
                          \item \al{花钱买}。
                          \item \al{将自己担任警察期间查办犯罪活动时掌握的甲抢劫财物的犯罪线索告诉检察人员，经查证属实}。
                          \item 犯罪分子将本人以往\zx{查办犯罪职务活动中}掌握的，或者从负有查办犯罪、监管职责的国家工作人员处获取的他人犯罪线索予以检举揭发的，不能认定为有立功表现。
                        \end{enumerate}
                  \item \zx{“新”指犯罪后形成}，不要求是否被公安掌握。
                        \begin{enumerate}
                          \item \al{供述体貌特征、供述手机号、家庭住址 \cuo{\xhx{犯罪前已客观存在}}}。
                          \item \al{供述\xhx{藏匿地址} \dui{犯罪后形成}}。
                        \end{enumerate}
                \end{enumerate}
          \item 协助抓捕。
          \item \zd{检举 xx 犯罪事实，构成立功}。
          \item 共同犯罪的自首立功。出卖兄弟三项好处。一个一般自首，两个立功（揭发犯罪、协助抓捕）。
        \end{enumerate}
  \item 数罪并罚。吸收、限制加重、并科、并罚、漏罪新罪。
        \begin{enumerate}
          \item 吸收原则。轻罪不执行了。死刑/无期 + 其他 = 死刑/无期、有期 + 拘役 = 有期。
          \item 限制加重原则。存在上限、下限，不能无限累加。考有期 + 有期。
                \begin{enumerate}
                  \item 下限。永远是受刑中的最高刑。\al{2, 5, 10 \yjt 10}。
                  \item 上限。从上往下，20、35 为界，最高 25。
                        \begin{enumerate}
                          \item \al{2, 5, 10。总刑期 \xyh 20 \yjt 总刑期 17, [10, 17]}。
                          \item \al{10, 11, 12。20 \xyh 总刑期 \xyh 35 \yjt 20, [12, 20]}。
                          \item \al{10, 12, 14。总刑期 \dy 35 \yjt 25, [14, 25]}。
                        \end{enumerate}
                \end{enumerate}
          \item 并科原则。有期 + 管制、拘役 + 管制分别执行。管制不关押，两者性质不同。
          \item 附加刑的并罚。种类相同合并执行。\al{剥夺政治权利 3 年、5 年}。不同分别执行。\al{罚金、没收财产}。
          \item 特殊情形。漏罪、新罪。\al{判了 5 年，执行了 3 年，新罪判 8 年}。
                \begin{enumerate}
                  \item 发现漏罪，先并后减。\al{[8, 13] 选一个 10 - 3 = 7}。
                  \item 又犯新罪，先减后并。\al{5 - 3 = 2。[8, 10] 选一个}。对被告更不利。
                  \item 都有。先漏后新。
                \end{enumerate}
        \end{enumerate}
  \item 缓刑、假释。参看刑诉。
\end{enumerate}
